\section*{第 \ref{ch:g11:seq} 章 --- 数列初步}

\begin{solution}{exo:g11:seq:1}
$u_n = \frac{n}{n+1}$: $u_1 = \frac12$, $u_2 = \frac23$,
$u_3 = \frac34$.

$u_0 = 5$, $u_{n+1} = 3u_n - 2$: $u_1 = 13$, $u_2 = 37$, $u_3 = 109$.

$u_n = (-1)^n n$: $u_1 = -1$, $u_2 = 2$, $u_3 = -3$.
\end{solution}

\begin{solution}{exo:g11:seq:2}
$u_{10} = 7 + 10 \times (-3) = -23$ 且 $u_{25} = 7 - 75 = -68$。

对于 $(v_n)$：由 $v_8 = v_3 + 5d$ 得 $26 = 11 + 5d$，所以 $d = 3$；
再得 $v_0 = v_3 - 3d = 11 - 9 = 2$。
\end{solution}

\begin{solution}{exo:g11:seq:3}
$u_8 = 5 \times 2^8 = 1280$.

对于 $(v_n)$：由 $v_4 = v_2\, q^2$ 得 $48 = 12 q^2$，故 $q^2 = 4$，且
$q = 2$（各项为正）。于是
$v_0 = \frac{v_2}{q^2} = \frac{12}{4} = 3$。
\end{solution}

\begin{solution}{exo:g11:seq:4}
$1 + \dots + 500 = \frac{500 \times 501}{2} = 125\,250$.

$4 + 7 + \dots + 61$ 为等差数列，公差 $d = 3$，且
共 $\frac{61 - 4}{3} + 1 = 20$ 项，其和为
$20 \times \frac{4 + 61}{2} = 650$.

$1 + \frac12 + \dots + \frac{1}{2^{10}}$ 为等比数列，公比 $q = \frac12$，
共 $11$ 项，其和为：
$\frac{1 - (1/2)^{11}}{1 - 1/2} = 2\left(1 - \frac{1}{2048}\right)
= \frac{2047}{1024}$.
\end{solution}

\begin{solution}{exo:g11:seq:5}
$u_{n+1} - u_n = 4(n+1) - 1 - 4n + 1 = 4$：为等差数列，公差 $d = 4$。

$\frac{v_{n+1}}{v_n} = \frac{2^{n+1}}{3^{n+2}} \cdot \frac{3^{n+1}}{2^n}
= \frac23$：为等比数列，公比 $q = \frac23$。

$w_0 = 0$，$w_1 = 2$，$w_2 = 6$：相邻两项之差 $2$ 与 $4$ 不相等，
故非等差数列；$\frac{w_1}{w_0}$ 无定义，且
比值 $\frac{w_2}{w_1} = 3 \neq \frac{w_3}{w_2} = 2$，故亦非等比数列。
\end{solution}

\begin{solution}{exo:g11:seq:6}
各排座位数成等差数列，首项 $16$，公差 $2$。最后
一排（第 20 排）有 $16 + 19 \times 2 = 54$ 个座位。
共 $20 \times \frac{16 + 54}{2} = 700$ 个座位。
\end{solution}

\begin{solution}{exo:g11:seq:7}
经过 $n$ 小时后，细菌数为 $500 \times 2^n$。晚上 8 点时，$n = 8$：
此时细菌有 $500 \times 256 = 128\,000$ 个。需满足 $500 \times 2^n >
10^6$，即 $2^n > 2000$：由于 $2^{10} = 1024$ 且 $2^{11} = 2048$，
细菌数首次超过一百万，是在经过 $11$ 个整小时后，即晚上 11 点。
\end{solution}

\begin{solution}{exo:g11:seq:8}
$c_2 = 1.002 \times 100 + 100 = 200.20$，且
$c_3 = 1.002 \times 200.20 + 100 \approx 300.60$。相邻两项之差
$c_2 - c_1 = 100.20$ 与 $c_3 - c_2 \approx 100.40$ 不相等，故
$(c_n)$ 不是等差数列；比值 $\frac{c_2}{c_1} = 2.002$ 与
$\frac{c_3}{c_2} \approx 1.50$ 亦不相等，故非等比数列。（此类“先乘后加”形式的
混合递推，可借辅助数列技巧求解，
见\cref{exo:g11:seq:11}。）
\end{solution}

\begin{solution}{exo:g11:seq:9}
$u_{n+1} - u_n = (n+1)^2 - 8(n+1) - n^2 + 8n = 2n - 7$：当
$n \leq 3$ 时为负，$n \geq 4$ 时为正。因此 $(u_n)$ 先递减，直至
$u_4 = 16 - 32 = -16$，之后递增，故该数列不单调。

$(v_n)$ 各项为正，且
\[
\frac{v_{n+1}}{v_n} = \frac{3^{n+1}}{(n+1)!} \cdot \frac{n!}{3^n}
= \frac{3}{n+1},
\]
该比值 $> 1$（当 $n \leq 1$ 时），$= 1$（$n = 2$ 时），且
$< 1$（$n \geq 3$ 时），故数列先递增至 $v_2 = v_3 = \frac92$，然后
递减。
\end{solution}

\begin{solution}{exo:g11:seq:10}
前 $n$ 项为 $u_0, \dots, u_{n-1}$，其中 $u_0 = 3$，且
$u_{n-1} = 3 + 4(n-1) = 4n - 1$。它们的和为
\[
n \times \frac{3 + (4n-1)}{2} = n(2n + 1) = 903,
\]
故 $2n^2 + n - 903 = 0$。此时 $\Delta = 1 + 4 \times 2 \times 903 =
7225 = 85^2$，且 $n = \frac{-1 + 85}{4} = 21$（负根
舍去）。检验：$21 \times 43 = 903$。
\end{solution}

\begin{solution}{exo:g11:seq:11}
\emph{1.} $u_1 = 5$，$u_2 = 9$，$u_3 = 17$：每一项都比
$4, 8, 16$ 多 1，由此推测 $u_n = 2^{n+1} + 1$。

\emph{2.} 令 $v_n = u_n - 1$：
\[
v_{n+1} = u_{n+1} - 1 = 2u_n - 1 - 1 = 2(u_n - 1) = 2v_n,
\]
故 $(v_n)$ 为等比数列，公比为 $2$，
首项为 $v_0 = u_0 - 1 = 2$。

\emph{3.} 于是 $v_n = 2 \times 2^n = 2^{n+1}$，
$u_n = v_n + 1 = 2^{n+1} + 1$，验证了猜想。（所减的数 $1$
在 $v_n$ 中，它是 $x \mapsto 2x - 1$ 的不动点；
同样的想法在 $u_{n+1} = au_n + b$ 的情形中再次出现，见\cref{ch:g12:seq}。）
\end{solution}

\begin{solution}{pb:g11:seq:1}
\textbf{1.} $h_1 = 1$, $h_2 = 3$, $h_3 = 7$.

\textbf{2.} 移动最大的圆盘时，其上方的 $n$ 个圆盘
必须先迁移到备用柱上（$h_n$ 次移动）；大圆盘
移过（$1$ 次移动）；随后这 $n$ 个圆盘必须重新
爬回其上（$h_n$ 次移动）：$h_{n+1} = 2h_n + 1$。

\textbf{3.} $v_{n+1} = h_{n+1} + 1 = 2h_n + 2 = 2v_n$:
公比为 $2$ 的等比数列，且 $v_1 = 2$，故 $v_n = 2^n$，从而
$h_n = 2^n - 1$.

\textbf{4.} $2^{64} - 1 \approx 1.8 \times 10^{19}$ 秒；
除以每年 $3 \times 10^7$ 秒，得
约 $6 \times 10^{11}$ 年 --- 即六千亿年，
约为宇宙年龄的四十倍。僧侣们
可以去喝咖啡休息。

\textbf{5.} 在任意合法解中，考虑最底下圆盘的
第一次移动：此时其余 $n$ 个圆盘必须全部位于
剩下的那根柱子上（至少需要 $h_n$ 次移动才能使它们到位），
而在最底下圆盘的最后一次移动之后，它们必须全部
重新回到其上（至少还需 $h_n$ 次移动）：任何解至少需要
$2h_n + 1$ 次移动。该递推式既是下界，
也是上界：$2^n - 1$ 是最优的。

\textbf{6.} $1, 1, 2, 3, 5, 8, 13, 21, 34, 55, 89, 144$.

\textbf{7.} 不是等差数列（$2 - 1 = 1$，但 $3 - 2 = 1$，
$5 - 3 = 2$：差值在变化）；不是等比数列
（$\frac21 = 2$，但 $\frac32 = 1.5$）。递增：当
$n \geq 2$ 时，$F_{n+1} - F_n = F_{n-1} > 0$。

\textbf{8.} $F_k = F_{k+2} - F_{k+1}$，于是
\[
\sum_{k=1}^{n} F_k = (F_3 - F_2) + (F_4 - F_3) + \dots +
(F_{n+2} - F_{n+1}) = F_{n+2} - F_2 = F_{n+2} - 1 .
\]
当 $n = 6$ 时：$1 + 1 + 2 + 3 + 5 + 8 = 20 = F_8 - 1 = 21 - 1$。

\textbf{9.} $F_k F_{k+1} - F_{k-1} F_k = F_k (F_{k+1} -
F_{k-1}) = F_k \cdot F_k = F_k^2$；求和裂项得
$F_n F_{n+1} - F_1 F_0$（其中 $F_0 = 0$）：平方和为
$F_n F_{n+1}$。当 $n = 4$ 时：$1 + 1 + 4 + 9 = 15 = F_4 F_5 =
3 \times 5$.

\textbf{10.} $F_5 F_3 - F_4^2 = 5 \times 2 - 9 = 1$;
$F_6 F_4 - F_5^2 = 8 \times 3 - 25 = -1$;
$F_7 F_5 - F_6^2 = 13 \times 5 - 64 = 1$：结果交替取
$\pm 1$。这一“差一”现象存在于 $F_{n+1} F_{n-1}$ 与
$F_n^2$ 之间，恰是魔术中“多出或少掉一个方格单位”的秘密：
将一个 $F_n \times F_n$ 的正方形剪开，重新拼成
$F_{n+1} \times F_{n-1}$ 的矩形时，必然多出或吞掉一个
单位 --- 即那条细缝。

\textbf{11.} $F_{n+2} = F_{n+1} + F_n \geq F_n + F_n = 2F_n$
（数列递增）：每隔两个下标，数值至少
翻一番 --- 相当于每个下标至少以公比 $\sqrt2$
的等比方式增长。

\textbf{12.} $1.5$; $1.667$; $1.6$; $1.625$; $1.615$;
$1.619$；$1.618$；$1.618$。若 $r_n \to L$，由
$F_{n+2} = F_{n+1} + F_n$，两边除以 $F_{n+1}$：
$r_{n+1} = 1 + \frac{1}{r_n}$，故 $L = 1 + \frac1L$，即\
$L^2 = L + 1$：$L = \varphi = \frac{1 + \sqrt5}{2}$，即
\cref{pb:g10:algebra:1} 中的黄金比。兔子便在
黄金中繁衍。

\textbf{13.} $v_{n+1} = u_{n+1} - \ell = a u_n + b - \ell$;
由 $\ell = a\ell + b$，上式化为 $a(u_n - \ell) = a v_n$：
公比为 $a$ 的等比数列。于是 $v_n = a^n v_0$，且
$u_n = a^n (u_0 - \ell) + \ell$.

\textbf{14.} 不动点：由 $\ell = 1.01\ell - 300$ 得
$\ell = 30\,000$。于是
$d_n = 1.01^n (10\,000 - 30\,000) + 30\,000
= 30\,000 - 20\,000 \times 1.01^n$.

\textbf{15.} 要使 $d_n \leq 0$，需 $1.01^n \geq 1.5$：
$1.01^{40} \approx 1.489$，$1.01^{41} \approx 1.504$：可见
第 $41$ 次还款即可还清债务（且略小于
$300$）。累计还款略低于 $41 \times 300 = 12\,300$
欧元 --- 所借的 $10\,000$ 欧元共付出约 $2\,300$ 欧元的
利息。

\textbf{16.} 不动点 $\ell = \frac{1000}{1 - 1.02} =
-50\,000$，故 $p_n = 1.02^n \times 100\,000 - 50\,000$。经过
$10$ 年后，$1.02^{10} \approx 1.219$：
$p_{10} \approx 71\,900$ 人。

\textbf{17.} $\frac{1000 \times 1001}{2} = 500\,500$；而
$2^{20} - 1 = 1\,048\,575$.

\textbf{18.} 从 $7$ 到 $502$，步长为 $5$：
$\frac{502 - 7}{5} + 1 = 100$ 项，其和
$= 100 \times \frac{7 + 502}{2} = 25\,450$.

\textbf{19.} 余额 $= 100 \times \frac{1.005^{60} - 1}
{1.005 - 1} \approx 100 \times \frac{0.3489}{0.005} \approx
6\,977$ 欧元 --- 其中存入 $6\,000$ 欧元，约 $977$
欧元为利息所得：等比数列求和正是银行的看家本领。

\textbf{20.} 显式公式可直接回答“$u_{1000}$ 是多少”；
递推关系则描述系统实际的演化方式 ---
其技巧在于把后者化为前者。等差数列
可以相加，等比数列可以相乘，每一族各有一个
求和公式（高斯的配对法；倍增法）。“不动点与辅助数列”的
技巧将每一个仿射
递推式化为等比数列 --- 贷款、人口与
汉诺塔都可用它解决。斐波那契数列不服从上述任一族，但
裂项恒等式仍刻画了它的和与平方和；其完整
面貌（精确通项公式与黄金比极限）有待更强的
工具。
\end{solution}
